% Part: first-order-logic
% Chapter: models-theories
% Section: set-theory

\documentclass[../../../include/open-logic-section]{subfiles}

\begin{document}

\olfileid{fol}{mat}{set}

\olsection{সেটের তত্ত্ব}

সেটের তত্ত্বে প্রায় সমস্ত গণিত গড়ে তোলা যায়। এই তত্ত্বে গণিত গড়ে
তোলার কয়েকটি দিক আছে। প্রথমত,~$\in$ সম্পর্কটির জন্য একটি
স্বতঃসিদ্ধসমষ্টি চাই। কখনও~$\in$-এর পরস্পরবিরোধী ধর্ম রেখেও নানা
স্বতঃসিদ্ধ-পদ্ধতি তৈরি হয়েছে। তবে নির্বাচন স্বতঃসিদ্ধসহ
জার্মেলো--ফ্রেঙ্কেল সেটতত্ত্ব নামে পরিচিত~$\Log{ZFC}$ বিশেষভাবে
উল্লেখযোগ্য। এটিই সবচেয়ে বেশি ব্যবহৃত ও অধ্যয়ন করা পদ্ধতি, কারণ দেখা
যায় গণিতবিদেরা যেসব ফল প্রমাণ করতে চান, তার প্রায় সবই প্রমাণ করার
জন্য এর স্বতঃসিদ্ধগুলি যথেষ্ট। কিন্তু তা প্রতিষ্ঠা করার আগে,
গণিতবিদেরা যা যা প্রকাশ করতে চান সেগুলি আদৌ কীভাবে প্রকাশ করা যায়,
তা স্পষ্ট করা দরকার। শুরুতেই সমস্যা হলো, ভাষাটিতে কোনো !!{constant}s
বা !!{function}s নেই। ফলে নির্দিষ্ট সেট (যেমন~$\emptyset$ বা~$\Nat$)
নিয়ে কথা বলা, সেটের উপর ক্রিয়া (যেমন~$X \cup Y$ ও~$\Pow{X}$) নিয়ে
কথা বলা, এমনকি সম্পর্ক ও অপেক্ষকের মতো সেট ছাড়া অন্য বস্তুযুক্ত
নির্মাণ নিয়ে কথা বলাও প্রথম দর্শনে অস্পষ্ট মনে হয়।

প্রথমে লক্ষ করি, “এর উপাদান” আমাদের আগ্রহের একমাত্র সম্পর্ক নয়;
“এর উপসেট” সম্পর্কটিও প্রায় সমান গুরুত্বপূর্ণ। কিন্তু “এর উপাদান”
সম্পর্কের সাহায্যে “এর উপসেট” সম্পর্কটি \emph{সংজ্ঞায়িত} করা যায়।
এর জন্য সেটতত্ত্বের ভাষায় এমন !!a{formula}~$!A(x,y)$ খুঁজতে হবে যা
সেটযুগল~$\tuple{X,Y}$ দ্বারা পরিতৃপ্ত হয় যদি এবং কেবল যদি
$X \subseteq Y$। কিন্তু~$X$ হলো~$Y$-এর উপসেট ঠিক তখনই, যখন~$X$-এর
সব !!{element}s,~$Y$-এরও !!{element}s। তাই নিচের সূত্র দিয়ে
$\subseteq$ সংজ্ঞায়িত করতে পারি:
\[
\lforall[z][(z \in x \lif z \in y)]
\]
এখন কোনো সূত্রে~$\subseteq$ সম্পর্কটি ব্যবহার করতে চাইলে, তার বদলে ওই
সূত্রটি বসাতে পারি ($x$ ও~$y$ যথাযথভাবে প্রতিস্থাপন করে এবং প্রয়োজন
হলে বদ্ধ চলরাশি~$z$-এর নাম বদলে)। যেমন, সেটের সদস্যভিত্তিক সমতার নীতি
বলে: যেকোনো দুটি সেট~$x$ ও~$y$ পরস্পরের মধ্যে অন্তর্ভুক্ত হলে
$x$ ও~$y$ একই সেট। একে
$\lforall[x][\lforall[y][((x \subseteq y \land y
    \subseteq x) \lif x = y)]]$ দিয়ে প্রকাশ করা যায়; অথবা উপরের
সংজ্ঞা দিয়ে~$\subseteq$ প্রতিস্থাপন করলে পাই
\[
\lforall[x][\lforall[y][((\lforall[z][(z \in x \lif z \in y)] \land
    \lforall[z][(z \in y \lif z \in x)]) \lif x = y)]].
\]
এটি আসলে~$\Log{ZFC}$-এর একটি স্বতঃসিদ্ধ, “সদস্যভিত্তিক সমতার
স্বতঃসিদ্ধ”।

$\emptyset$-এর জন্য কোনো !!{constant} নেই, কিন্তু
$\lnot \lexists[y][y \in x]$ দিয়ে “$x$ শূন্য” প্রকাশ করা যায়। তখন
“$\emptyset$ আছে” হয়ে যায়
!!{sentence}~$\lexists[x][\lnot \lexists[y][y \in x]]$। এটি
$\Log{ZFC}$-এর আরেকটি স্বতঃসিদ্ধ। (লক্ষ করো, সদস্যভিত্তিক সমতার
স্বতঃসিদ্ধ থেকে অনুসৃত হয় যে শূন্য সেট একটিই।) সেটতত্ত্বের ভাষায়
$\emptyset$ নিয়ে কথা বলতে চাইলে লিখতে হবে, “এমন একটি সেট আছে যা শূন্য
এবং \dots”। যেমন, $\emptyset$ প্রত্যেক সেটের উপসেট—এ কথা প্রকাশ করতে
লিখতে পারি
\[
\lexists[x][(\lnot\lexists[y][y \in x] \land \lforall[z][x \subseteq
    z])]
\]
যেখানে অবশ্য $x \subseteq z$-কেও আবার তার সংজ্ঞা দিয়ে প্রতিস্থাপন
করতে হবে।

$X \cup Y$ ও~$\Pow{X}$-এর মতো সেটের উপর ক্রিয়া নিয়ে কথা বলতে একই
কৌশল নিতে হয়। সেটতত্ত্বের ভাষায় কোনো অপেক্ষক-সংকেত নেই, কিন্তু নিচের
সূত্রগুলি দিয়ে অপেক্ষকধর্মী সম্পর্ক~$X \cup Y = Z$ ও~$\Pow{X} = Y$
প্রকাশ করা যায়:
\begin{align*}
& \lforall[u][((u \in x \lor u \in y) \liff u \in z)]\\
& \lforall[u][(u \subseteq x \liff u \in y )]
\end{align*}
কারণ $X \cup Y$-এর !!{element}s ঠিক সেই সেটগুলি যেগুলি~$X$-এর
!!{element}s অথবা~$Y$-এর !!{element}s, এবং~$\Pow{X}$-এর
!!{element}s ঠিক~$X$-এর উপসেটগুলি।
কিন্তু এর ফলে $x \cup y$ বা~$\Pow{x}$-কে পদের মতো ব্যবহার করা যায় না;
শুধু~$X \cup Y = Z$ ও~$\Pow{X} = Y$ সম্পর্ক দুটিকে সংজ্ঞায়িত করা
সম্পূর্ণ !!{formula}s ব্যবহার করতে পারি। এমনকি এই সম্পর্কগুলি কখনও
পরিতৃপ্ত হয় কি না, অর্থাৎ সংঘ ও ঘাত সেট সব সময় আছে কি না, তা এখনও
জানা নেই। যেমন, !!{sentence}
$\lforall[x][\lexists[y][\Pow{x} = y]]$,~$\Log{ZFC}$-এর আরেকটি
স্বতঃসিদ্ধ (ঘাত সেটের স্বতঃসিদ্ধ)।

এবার ক্রমযুগল বা অপেক্ষক নিয়ে কীভাবে কথা বলব? এখানে ক্রমযুগল ও
অপেক্ষককে বিশেষ ধরনের সেট হিসেবে কীভাবে ভাবা যায়, তা ব্যাখ্যা করতে
হবে। ক্রমযুগল~$\tuple{x,y}$-কে সংজ্ঞায়িত করার একটি উপায় হলো
$\{\{x\},\{x,y\}\}$ সেটটি। কিন্তু আগের মতো এখানেও এই সেটটির নাম দেয়
এমন কোনো !!a{function} প্রবর্তন করা যায় না; শুধু
$\tuple{x,y}=z$, অর্থাৎ $\{\{x\},\{x,y\}\}=z$ সম্পর্কটি সংজ্ঞায়িত
করা যায়:
\[
\lforall[u][(u \in z \liff (\lforall[v][(v \in u \liff v = x)] \lor
  \lforall[v][(v \in u \liff (v = x \lor v = y))]))]
\]
এটি বলে,~$z$-এর !!{element}s~$u$ ঠিক সেই সেটগুলি যাদের একমাত্র
!!{element}~$x$, অথবা যাদের একমাত্র !!{element}s~$x$ ও~$y$ (অন্যভাবে
বললে, ঠিক সেই সেটগুলি যেগুলি হয়~$\{x\}$-এর সঙ্গে, নয়তো
$\{x,y\}$-এর সঙ্গে অভিন্ন)। এটি পেয়ে গেলে আরও কথা বলা যায়; যেমন,
$X \times Y = Z$:
\[
\lforall[z][(z \in Z \liff \lexists[x][\lexists[y][(x \in
        X \land y \in Y \land \tuple{x, y} = z)]])]
\]

কোনো অপেক্ষক~$f \colon X \to Y$-কে $f(x)=y$ সম্পর্ক, অর্থাৎ
যুগলসমষ্টি~$\Setabs{\tuple{x,y}}{f(x)=y}$ হিসেবে ভাবা যায়। তখন বলতে
পারি, কোনো সেট~$f$,~$X$ থেকে~$Y$-তে একটি অপেক্ষক, যদি (ক) সেটি একটি
সম্পর্ক~$\subseteq X \times Y$, (খ) সেটি সর্বত্র সংজ্ঞায়িত, অর্থাৎ
প্রত্যেক $x \in X$-এর জন্য এমন কোনো $y \in Y$ আছে যাতে
$\tuple{x,y} \in f$, এবং (গ) সেটি অপেক্ষকধর্মী, অর্থাৎ যখনই
$\tuple{x,y},\tuple{x,y'} \in f$, তখন $y=y'$ (কারণ অপেক্ষকের মান
অনন্য হতে হয়)। তাই “$f$,~$X$ থেকে~$Y$-তে একটি অপেক্ষক” এভাবে লেখা
যায়:
\begin{align*}
 & \lforall[u][(u \in f \lif
      \lexists[x][\lexists[y][(x \in X \land y \in
      Y \land \tuple{x, y} = u)]])] \land {}\\
 & \lforall[x][(x \in X \lif
      (\lexists[y][(y \in Y \land \mathrm{maps}(f, x, y))] \land
      \lforall[y][\lforall[y'][((\mathrm{maps}(f, x, y) \land
      \mathrm{maps}(f, x, y')) \lif y = y')]]))]
\end{align*}
% BN-SRC-105 TARGET-CORRECTION-BEGIN
\par\smallskip\noindent\textnormal{\small\textbf{উৎস-সংশোধন (BN-SRC-105):} হিমায়িত ইংরেজি প্রদর্শনে উভয় সার্বিক পরিমাণসূচকের সূত্র-আর্গুমেন্টই পূর্বপক্ষের পরে অকালে বন্ধ হয়েছে; ফলে অস্তিত্বমূলক শর্তগুলি শর্তবাচকের পরপক্ষের বাইরে পড়ে এবং দ্বিতীয় সারির বন্ধনীগুলিও অসম্পূর্ণ থাকে। বাংলা পাঠে ঘোষিত তিনটি শর্ত মেনে প্রথম পরপক্ষে যুগল-সদস্যতার অস্তিত্ব এবং দ্বিতীয় পরপক্ষে অস্তিত্ব ও মানের অনন্যতা রাখা হয়েছে, সঙ্গে সব পরিমাণসূচক ও বন্ধনী সামঞ্জস্য করা হয়েছে।} % BN-SRC-105 TARGET-CORRECTION-END
যেখানে $\mathrm{maps}(f,x,y)$ হলো
$\lexists[v][(v \in f \land \tuple{x,y}=v)]$-এর সংক্ষিপ্ত রূপ
(এই !!{formula} “$f(x)=y$” প্রকাশ করে)।

এখন, যেমন $f\colon X\to Y$ !!{injective}—এ কথাও প্রকাশ করা কঠিন নয়:
\begin{multline*}
 f \colon X \to Y \land \lforall[x][\lforall[x'][((x \in X \land x' \in
     X \land {} \\
 \lexists[y][(\mathrm{maps}(f, x, y) \land \mathrm{maps}(f,
        x', y))]) \lif x = x')]]
\end{multline*}
% BN-SRC-106 TARGET-CORRECTION-BEGIN
\par\smallskip\noindent\textnormal{\small\textbf{উৎস-সংশোধন (BN-SRC-106):} হিমায়িত ইংরেজি injectivity-সূত্রে প্রথম সারির শেষে দুটি অকাল সমাপ্তি-বর্গবন্ধনী আছে, অথচ দ্বিতীয় সারির শেষে দুই সার্বিক পরিমাণসূচকের সমাপ্তি-বর্গবন্ধনী অনুপস্থিত। বাংলা পাঠে যৌথ পূর্বপক্ষের ভিতরে একই মানে পাঠানোর অস্তিত্বমূলক সূত্রটি রেখে, তারপর শর্তবাচক ও দুই সার্বিক পরিমাণসূচক যথাস্থানে বন্ধ করা হয়েছে।} % BN-SRC-106 TARGET-CORRECTION-END
কোনো অপেক্ষক~$f\colon X\to Y$ !!{injective} যদি এবং কেবল যদি
$f$,~$x,x'\in X$-কে একই~$y$-তে পাঠালেই $x=x'$ হয়। এই সূত্রটিকে
$\mathrm{inj}(f,X,Y)$ দিয়ে সংক্ষেপ করলে সেটতত্ত্বের ভাষায় কান্টরের
উপপাদ্যের মতো অ-তুচ্ছ কথাও ইতিমধ্যে বলা যায়:~$\Pow{X}$ থেকে~$X$-এ
কোনো !!{injective} অপেক্ষক নেই:
\[
\lforall[X][\lforall[Y][(\Pow{X} = Y \lif
    \lnot\lexists[f][\mathrm{inj}(f, Y, X)])]]
\]

মনে হতে পারে, প্রত্যেক সংজ্ঞায়ক ধর্মের জন্য একটি সেটের অস্তিত্ব নিশ্চিত
করতে সেটতত্ত্বে আরেকটি স্বতঃসিদ্ধ দরকার। $!A(x)$ যদি মুক্ত
চলরাশি~$x$-সহ সেটতত্ত্বের একটি সূত্র হয়, তবে
!!{sentence}
\[
\lexists[y][\lforall[x][(x \in y \liff !A(x))]].
\]
বিবেচনা করতে পারি। এই !!{sentence} বলে, এমন একটি সেট~$y$ আছে যার
!!{element}s ঠিক সেই~$x$-গুলি যা~$!A(x)$ পরিতৃপ্ত করে। এই ছকটির নাম
“ধর্মনির্দেশে সেট-গঠন নীতি”। এটি খুব উপকারী মনে হয়; দুর্ভাগ্যবশত এটি
অসঙ্গত। $!A(x) \ident \lnot x \in x$ নিলে নীতিটি বলে
\[
\lexists[y][\lforall[x][(x \in y \liff x \notin x)]],
\]
অর্থাৎ, নিজের !!{element}s নয় এমন সব সেটের একটি সেট আছে। এমন কোনো সেট
থাকতে পারে না—এটিই রাসেলের কূটাভাস। আসলে~$\Log{ZFC}$-তে এই নীতির একটি
সীমিত—এবং সঙ্গত—রূপ, পৃথকীকরণ নীতি, আছে:
\[
\lforall[z][\lexists[y][\lforall[x][(x \in y \liff (x \in z \land
      !A(x)))]]].
\]
% BN-SRC-107 TARGET-CORRECTION-BEGIN
\par\smallskip\noindent\textnormal{\small\textbf{উৎস-সংশোধন (BN-SRC-107):} হিমায়িত ইংরেজি পৃথকীকরণ-নীতির সূত্রে অন্তঃস্থ সংযোজনী বন্ধ করার পরে biconditional-এর উদ্বোধনী গোলবন্ধনীটি বন্ধ করা হয়নি। বাংলা পাঠে তিনটি সমাপ্তি-বর্গবন্ধনীর আগে সেই সমাপ্তি-গোলবন্ধনী পুনরুদ্ধার করা হয়েছে।} % BN-SRC-107 TARGET-CORRECTION-END

\begin{prob}
!!a{derivation} দিয়ে
\[
\lexists[y][\lforall[x][(x \in y \liff x \notin x)]] \Proves \lfalse.
\]
দেখিয়ে প্রমাণ করো যে ধর্মনির্দেশে সেট-গঠন নীতি অসঙ্গত।
আগে $(A \lif \lnot A) \land (\lnot A \lif A)
\Proves \lfalse$ দেখানো সহায়ক হতে পারে।
\end{prob}
\end{document}
